%%SECTION abstract
For a person and a language model reading the same text, a simple linear map predicts the person's brain activity from the model's middle-layer activations. \evd{C1}{supported}
For a decade this correspondence has been used in one direction only: to measure a model's alignment with the brain after it is trained.
We ask whether it can instead be used during training.
We test this in the setting of model compression: we distill a language model into a smaller student with brain alignment in the objective, at a matched budget, and ask whether the student preserves or improves that alignment, a question that remains untested.
An information-budget argument implies that such a signal is far too small to act as a second teacher; at most it is a weak prior.
Such a prior should help only where the text signal is weakest (small models, scarce data, aggressive compression) and wash out wherever data is abundant.
\gap{Results summary deferred (checkpoint 2): the two established findings and the downstream arc are summarised here only once R08--R14 are consolidated.}
